\documentclass[11pt,a4paper]{article}
\usepackage[margin=25mm]{geometry}
\usepackage{amsmath,amssymb,amsthm,hyperref}
\setlength{\emergencystretch}{2em}
\newtheorem{lemma}{Lemma}
\newtheorem{theorem}[lemma]{Theorem}
\newcommand{\gen}[1]{\langle#1\rangle}
\title{Excluding endpoint touching bands}
\author{Human Compression project: L044 continuation}
\date{}
\begin{document}\maketitle
Use the standing hypotheses of the accepted nested-bridge theorem:
$G\leq A_n$, $n\geq12$, is transitive with a connected string
C-marking and all maximal omissions intransitive. Suppose its
high-rank equality $g=h$ occurs, with one crossing at $i=0$
(the other orientation is the literal reversal). The unique edges
are $a_0\mathrel{-_0}b_0$ and $c\mathrel{-_j}d$, cutting the full
graph into $A,M,D$. Thus
\[
 n=2r,\quad |A|=h+1,\quad |M|=j,\quad
 |D|=2r-j-h-1,\quad 0<h<j.
\]
The right suffix $G_{>j}$ fixes $A\cup M$.
Retain the exact last/first shared-colour hypotheses: $h$ is
last shared across the $0$-cut, and $g=h$ first shared across
the $j$-cut. Consequently
\[
 \rho_0|_D=1,\qquad \rho_j|_A=1,
 \qquad G_{>h}\text{ fixes }A.                         \tag{1}
\]
All claims below keep these genuine purity conditions.

\begin{lemma}[The endpoint forces a three-point outside orbit]
Under these assumptions $h=2$, and, writing the forced ordered
paths as below,
\[
                  \rho_0=(a_0\ b_0)(a_1\ a_2).         \tag{2}
\]
\end{lemma}
\begin{proof}
The accepted fixed-prefix path in $A$, including its last
$h$-edge, has at least $h+1$ vertices and therefore exhausts $A$.
The middle path between the bridges similarly exhausts $M$.
Each intermediate colour occurs exactly once, and their order
is forced from the first anchor: all colours at least two fix
that anchor; after following colour one all colours at least
three fix the next point, by commutation; continuing gives
\[
 A:\ a_0\mathrel{-_1}a_1\mathrel{-_2}\cdots
       \mathrel{-_h}a_h,\qquad
 M:\ b_0\mathrel{-_1}b_1\mathrel{-_2}\cdots
       \mathrel{-_{j-1}}b_{j-1}=c.
\]
The suffix of the middle path from $b_1$ to $c$, followed by
the $j$-edge to $d$, uses only indices at least two. The element
$\rho_0$ fixes $d$ by (1) and commutes with every generator on
that path. Hence it fixes $M\setminus\{b_0\}$ pointwise, as
well as all of $D$. Its unique cross-orbit transposition is
$(a_0\ b_0)$, so its remaining action is inside
$A'=\{a_1,\ldots,a_h\}$.

For $h=1$ that remaining set is a singleton, making $\rho_0$
odd. For $h\geq3$, the group $\gen{\rho_2,\ldots,\rho_h}$
is transitive on $A'$ by the displayed path. Its subgroup
$\gen{\rho_3,\ldots,\rho_h}$ has exactly one fixed point in
$A'$, namely $a_1$: all other points are moved by displayed
edges of those colours, while $a_1$ is fixed by commutation
from the anchor. The restriction of $\rho_0$ centralizes both
groups, so fixes $a_1$ and then all of $A'$. This again makes
it odd. Thus $h=2$, and evenness forces (2).
\end{proof}

\begin{theorem}[No endpoint touching equality]
The high-rank equality $g=h$ is impossible when at least one
crossing is an endpoint and the stated first/last shared-colour
and orientation hypotheses hold.
\end{theorem}
\begin{proof}
If $j=r-1$, apply the lemma in both index directions. It gives
$h=2$ and $j-h=2$, so $r=5$ and $n=10$, contrary to $n\geq12$.

Suppose $j<r-1$. Strictness in the reversed outside-orbit bound
forces $h=j-1$. Thus $h=2$, $j=3$. Put $l=r-j=r-3\geq2$.
Then $|D|=2l$, $M=\{b_0,b_1,c\}$, and on $A\cup M$ the
first three generators are exactly
\[
 \rho_0=(a_0\ b_0)(a_1\ a_2),\quad
 \rho_1=(a_0\ a_1)(b_0\ b_1),\quad
 \rho_2=(a_1\ a_2)(b_1\ c)\delta,                    \tag{3}
\]
where $\delta=\rho_2|_D$. Indeed $\rho_1$ fixes $D$ because
$g=2$ is the first shared colour, and it fixes $c$ by the
unique $3$-crossing; the remaining actions in (3) follow
from the three-point paths. Write $R=G_{>3}$.
The group $\gen{\delta,R}$ is transitive on $D$, and $\delta$
centralizes $R$. The same critical commuting-orbit argument as
in the accepted interior note applies on this side alone:
$\delta$ is fixed-point-free, exchanging $l$ pairs, and either
$R$ is transitive on $D$ or it has two faithful natural $S_l$
orbits of size $l$. To spell out the count, a transitive $R$
has essential independent pair images of length $l-1$, hence
full $S_l$ image by Whiston; an intransitive $R$ has faithful
orbits of size at least $l$, hence exactly two, again $S_l$.
This uses only the actual global omission argument and not
an intersection property of a projection.

Equation (3) and evenness make $l$ even. Since $n=2l+6\geq12$,
it follows that $l\geq4$. Set $t=\rho_3$ and
$Q=\gen{t,R}=G_{>2}$. The element $t$ fixes $A$ by (1), fixes $b_0$ by the unique
$0$-crossing, and fixes $b_1=b_0^{\rho_1}$ by commutation with
$\rho_1$. Hence $t=(c\ d)v$, where $v$ is supported on
$D\setminus\{d\}$. We next show that $Q$
is transitive on $Z=D\cup\{c\}$.
If $R$ is transitive on $D$, the $cd$-edge proves this immediately.
Otherwise identify the two $S_l$ rows through $\delta$.
The subgroup $H=G_{>4}$ fixes $d,d'=d^\delta$. Its pair
images form an independent family of length $l-2$ in $S_{l-1}$,
so Whiston's equality theorem makes its action on the remaining
$l-1$ pairs the full $S_{l-1}$. In this two-row case it is the
same natural action in both rows. The involution $v$ centralizes
$H$ and fixes $d'$: it preserves $D\setminus\{d\}$, whose only
$H$-fixed point is $d'$. On the two remaining rows of length
$l-1\geq3$, a commuting permutation is either the identity
or the exchange identifying those rows, because the natural
$S_{l-1}$ centralizer is trivial. Since $t$ is even, $v$ must
be odd and therefore is that row exchange. Thus $t$ joins the
two $R$-orbits as well as adjoining $c$, proving transitivity.

The actual supported C-group $Q\leq A_Z$ has degree $2l+1$
and rank $l$, with the anchor $c$ fixed by $R$. Every maximal
omission inside $Q$ is intransitive on $Z$: otherwise the
corresponding global omission is already transitive on $Z$,
and adjoining its discarded generator, supported on $Z$, cannot
join omission orbits, contrary to transitivity of $G$.
If $|Z|=9$, the complete accepted small supported-tail lemma
already excludes $Q$: at this rank its only small survivors
have degrees five and six. That lemma's proof uses the actual
extension anchor, the complete small primitive attachment
certificate and the small anchored-imprimitive bound; it uses
no smaller-degree alternating induction. Thus we may assume
$|Z|\geq13$. The accepted proper-primitive bound and
anchored-imprimitive bound then both exclude a proper transitive
$Q<A_Z$ at this rank. Consequently
\[
                              Q=A_Z.                    \tag{4}
\]
All required small census programs and classical premises of
those exact predecessors remain paid; no new census is invoked.

The six-point core in (3) is exactly the accepted Table~2(2)
$A_5$ marking. Let $P=\gen{\rho_0,\rho_1,\rho_2}$.
It is a subgroup of $A_5\times\gen\delta$, surjecting onto
both factors. If $\delta\notin P$, it would be the graph of
a nontrivial homomorphism $A_5\to C_2$, impossible by
perfection of $A_5$. Thus $\delta\in P$.
But $\delta$ is even, nontrivial and supported on $D$, so
(4) puts it in $Q$ too. Full IP gives $P\cap Q=1$ because
their distinguished index sets are disjoint. This is the final
contradiction. Reversing handles a right endpoint crossing.
\end{proof}

\paragraph{Attribution and exact boundary.}
This is a new replacement argument for the endpoint touching
case left by the two-bridge count in Cameron--Fernandes--Leemans--Mixer,
\emph{Highest rank of a polytope for $A_n$}, arXiv:1605.09173v1,
Proposition~5.28, physical page~37. All accepted path, outside-orbit,
critical-pair, proper-primitive, anchored and small-tail proofs
are charged as whole dependencies, including their finite
programs and retained classical inputs. The proof does not
reclassify the source tables or repeat those computations.

Together with the accepted interior touching argument, it forces
$g>h$ in the oriented high-rank bridge situation. The subsequent
$g>h$ support and middle-orbit equality analysis remains separate.
There is no whole L044, new formal theorem, original, question
or priority credit in this bounded packet.
\end{document}
